\subsection*{CU-04: Cambiar la contraseña}
\addcontentsline{toc}{subsection}{CU-04: Cambiar la contraseña}
\label{cu-04}

\begin{fichacu}{CU-04}{Cambiar la contraseña}
  \crudFila{Responsable}{José Eduardo Prior Hernández}
  \crudFila{Actor principal}{Jugador}
  \crudFila{Actores de sistema}{Servidor de partidas, Servicio de correo, Base de datos}
  \crudFila{Prototipos}{9f}
  \crudFila{Prioridad}{Must (debe estar en la entrega). Categoría (a) Obligatorio: característica 2 del cliente (cambio de contraseña).}
  \crudFila{Objetivo}{Permitir al Jugador con sesión abierta cambiar su contraseña, exigiendo la contraseña actual como prueba de identidad y cerrando las demás sesiones de la cuenta.}
  \crudFila{Disparador}{El Jugador selecciona ``cambiar'' junto a la contraseña en la \texttt{GUI\_\allowbreak{}AccountSettings}.}
  \textbf{Precondiciones} & \cuCelda{\begin{cureglas}
      \item \textbf{\mbox{PRE-01}} El Jugador tiene una \textit{Sesion} abierta y su token es válido.
      \item \textbf{\mbox{PRE-02}} El \textit{Usuario} tiene \texttt{tipo\_\allowbreak{}cuenta} en \texttt{REGISTRADA} y \texttt{estado\_\allowbreak{}cuenta} en \texttt{ACTIVA}.
      \item \textbf{\mbox{PRE-03}} El Servidor de partidas está en ejecución y tiene conexión disponible con la Base de datos.
    \end{cureglas}} \\ \hline
  \textbf{Flujo normal} & \cuCelda{\begin{cuenum}[start=1]
      \item El Jugador selecciona ``cambiar'' junto a la contraseña.
      \item El SJM despliega la \texttt{GUI\_\allowbreak{}ChangePassword}, solicitando la contraseña actual, la contraseña nueva y su confirmación, con las opciones ``Guardar'' y ``Cancelar''. \textit{(Ver \mbox{FA-01})}
      \item El Jugador captura los tres campos y da click en ``Guardar''.
      \item El SJM valida que los campos estén completos, que la contraseña nueva cumpla la política de \mbox{RN-03} y que coincida con su confirmación. \textit{(Ver \mbox{FA-02})}
      \item El SJM envía el mensaje \texttt{CHANGE\_\allowbreak{}PASSWORD\_\allowbreak{}REQUEST} con la contraseña actual, la nueva y el token de sesión. \textit{(Ver \mbox{EX-03}, \mbox{EX-04}, \mbox{EX-05})}
      \item El Servidor de partidas verifica que el token corresponda a una \textit{Sesion} abierta. \textit{(Ver \mbox{FA-03})}
      \item El Servidor de partidas deriva el resumen de la contraseña actual recibida con la \texttt{contrasena\_\allowbreak{}sal} del \textit{Usuario} y lo compara contra \texttt{contrasena\_\allowbreak{}hash} (\mbox{RN-01}). \textit{(Ver \mbox{FA-04})}
      \item El Servidor de partidas vuelve a validar la política de la contraseña nueva y comprueba que no coincida con la actual (\mbox{RN-06}). \textit{(Ver \mbox{FA-05})}
    \end{cuenum}} \\
   & \cuCelda{\begin{cuenum}[start=9]
      \item El Servidor de partidas inicia una transacción sobre la Base de datos.
      \item El Servidor de partidas genera una \texttt{contrasena\_\allowbreak{}sal} nueva y deriva la \texttt{contrasena\_\allowbreak{}hash} de la contraseña nueva (\mbox{RN-02}). \textit{(Ver \mbox{EX-01})}
      \item El Servidor de partidas cierra todas las \textit{Sesion} abiertas del \textit{Usuario} \textbf{salvo la que envió la solicitud}, con \texttt{motivo\_\allowbreak{}cierre} igual a \texttt{CAMBIO\_\allowbreak{}CREDENCIALES} (\mbox{RN-04}).
      \item El Servidor de partidas invalida los \textit{CodigoVerificacion} pendientes del \textit{Usuario}, como el código de prueba de una activación del segundo factor sin terminar (\mbox{CU-05}), porque pudo solicitarlos quien conocía la contraseña anterior.
      \item El Servidor de partidas confirma la transacción. \textit{(Ver \mbox{EX-02})}
      \item El Servidor de partidas registra el cambio en la \textit{BitacoraAcceso} con el resultado \texttt{CAMBIO\_\allowbreak{}CONTRASENA}.
      \item El Servidor de partidas solicita al Servicio de correo un aviso al \texttt{correo} del \textit{Usuario} informando del cambio, en su \texttt{idioma\_\allowbreak{}preferido} y fuera de la transacción (\mbox{RN-05}). \textit{(Ver \mbox{EX-06})}
      \item El Servidor de partidas responde con \texttt{CHANGE\_\allowbreak{}PASSWORD\_\allowbreak{}OK} y el número de sesiones cerradas.
    \end{cuenum}} \\
   & \cuCelda{\begin{cuenum}[start=17]
      \item El SJM muestra la \texttt{GUI\_\allowbreak{}MessageSuccess} indicando que la contraseña se cambió y en cuántos dispositivos se cerró la sesión, y vuelve a la \texttt{GUI\_\allowbreak{}AccountSettings}.
      \item Fin del caso de uso.
    \end{cuenum}} \\ \hline
  \textbf{Flujos alternos} & \cuCelda{\textbf{\mbox{FA-01}: El Jugador cancela el cambio}\par
    \begin{cuenum}[start=1]
      \item El Jugador selecciona ``Cancelar''.
      \item El SJM descarta lo capturado sin enviarlo y vuelve a la \texttt{GUI\_\allowbreak{}AccountSettings}.
      \item Fin del caso de uso.
    \end{cuenum}} \\
   & \cuCelda{\textbf{\mbox{FA-02}: Campos incompletos, contraseña fuera de política o que no coincide}\par
    \begin{cuenum}[start=1]
      \item El SJM resalta los campos afectados y escribe junto a cada uno qué se espera, conforme a \mbox{RN-03}.
      \item El SJM limpia los campos de contraseña.
      \item El flujo continúa en el paso 3 del FN. La validación es local y no llega al Servidor de partidas.
    \end{cuenum}} \\
   & \cuCelda{\textbf{\mbox{FA-03}: La sesión ya no es válida}\par
    \begin{cuenum}[start=1]
      \item El Servidor de partidas responde con \texttt{SESSION\_\allowbreak{}INVALID}.
      \item El SJM muestra la \texttt{GUI\_\allowbreak{}MessageWarning} indicando que la sesión terminó y despliega la \texttt{GUI\_\allowbreak{}Login}.
      \item Fin del caso de uso.
    \end{cuenum}} \\
   & \cuCelda{\textbf{\mbox{FA-04}: La contraseña actual es incorrecta}\par
    \begin{cuenum}[start=1]
      \item El Servidor de partidas incrementa \texttt{intentos\_\allowbreak{}fallidos} y actualiza \texttt{fecha\_\allowbreak{}ultimo\_\allowbreak{}intento\_\allowbreak{}fallido} del \textit{Usuario}, aplicando el bloqueo escalonado del \mbox{RN-03} del \mbox{CU-01} (\mbox{RN-07}).
      \item El Servidor de partidas registra el intento en la \textit{BitacoraAcceso} con el resultado \texttt{CONTRASENA\_\allowbreak{}INCORRECTA}.
      \item Si la cuenta quedó bloqueada, el Servidor de partidas cierra todas sus \textit{Sesion} con \texttt{motivo\_\allowbreak{}cierre} igual a \texttt{CAMBIO\_\allowbreak{}CREDENCIALES} y responde con \texttt{ACCOUNT\_\allowbreak{}LOCKED}; el SJM despliega la \texttt{GUI\_\allowbreak{}Login} con el aviso de bloqueo y termina el caso de uso.
      \item En caso contrario, el Servidor de partidas responde con \texttt{CHANGE\_\allowbreak{}CREDENTIALS\_\allowbreak{}FAILED} y el motivo \texttt{CONTRASENA\_\allowbreak{}ACTUAL\_\allowbreak{}INCORRECTA}; el SJM limpia el campo y el flujo continúa en el paso 3 del FN.
    \end{cuenum}} \\
   & \cuCelda{\textbf{\mbox{FA-05}: La contraseña nueva es igual a la actual}\par
    \begin{cuenum}[start=1]
      \item El Servidor de partidas responde con \texttt{CHANGE\_\allowbreak{}CREDENTIALS\_\allowbreak{}FAILED} y el motivo \texttt{MISMA\_\allowbreak{}CONTRASENA}.
      \item El SJM muestra la \texttt{GUI\_\allowbreak{}MessageWarning} indicando que debe elegir una distinta y limpia los campos.
      \item El flujo continúa en el paso 3 del FN.
    \end{cuenum}} \\ \hline
  \textbf{Excepciones} & \cuCelda{\textbf{\mbox{EX-01}: Fallo al escribir en la base de datos}\par
    \begin{cuenum}[start=1]
      \item El Servidor de partidas revierte la transacción, de modo que no quede la contraseña cambiada con las demás sesiones todavía abiertas.
      \item El Servidor de partidas registra la excepción en la bitácora del servidor con un identificador de incidencia, sin incluir ninguna contraseña recibida.
      \item El Servidor de partidas responde con \texttt{SERVER\_\allowbreak{}ERROR} y el identificador.
      \item El SJM muestra la \texttt{GUI\_\allowbreak{}MessageError} con el identificador y limpia los campos de contraseña.
      \item El flujo continúa en el paso 3 del FN.
    \end{cuenum}} \\
   & \cuCelda{\textbf{\mbox{EX-02}: Fallo al confirmar la transacción}\par
    \begin{cuenum}[start=1]
      \item El Servidor de partidas trata la transacción como fallida y registra la excepción señalando que el estado es indeterminado.
      \item El Servidor de partidas responde con \texttt{SERVER\_\allowbreak{}ERROR}.
      \item El SJM muestra la \texttt{GUI\_\allowbreak{}MessageError} indicando que el cambio pudo haberse aplicado y que, si la contraseña nueva no funciona al volver a entrar, use la anterior.
      \item Fin del caso de uso.
    \end{cuenum}} \\
   & \cuCelda{\textbf{\mbox{EX-03}: Se agota el tiempo de espera de la respuesta}\par
    \begin{cuenum}[start=1]
      \item El SJM detecta que transcurrió el tiempo límite sin recibir un mensaje completo.
      \item El SJM no reenvía la solicitud y muestra la \texttt{GUI\_\allowbreak{}MessageError} indicando que el cambio pudo haberse aplicado.
      \item Fin del caso de uso.
    \end{cuenum}} \\
   & \cuCelda{\textbf{\mbox{EX-04}: La conexión se interrumpe durante el intercambio}\par
    \begin{cuenum}[start=1]
      \item El SJM detecta el cierre inesperado del socket antes de recibir respuesta.
      \item El flujo continúa en el paso 2 del \mbox{EX-03}.
    \end{cuenum}} \\
   & \cuCelda{\textbf{\mbox{EX-05}: El mensaje recibido está malformado}\par
    \begin{cuenum}[start=1]
      \item El SJM descarta el mensaje, cierra la conexión y registra en la bitácora local el contenido truncado.
      \item El SJM muestra la \texttt{GUI\_\allowbreak{}MessageError} indicando que la respuesta no pudo interpretarse.
      \item Fin del caso de uso.
    \end{cuenum}} \\
   & \cuCelda{\textbf{\mbox{EX-06}: El servicio de correo no está disponible}\par
    \begin{cuenum}[start=1]
      \item El Servidor de partidas agota los reintentos previstos hacia el Servicio de correo.
      \item El Servidor de partidas conserva el cambio de contraseña, que ya está confirmado, y registra el fallo del aviso en la bitácora del servidor (\mbox{RN-05}).
      \item El Servidor de partidas responde con \texttt{CHANGE\_\allowbreak{}PASSWORD\_\allowbreak{}OK}, indicando que el aviso no pudo enviarse.
      \item El flujo continúa en el paso 17 del FN.
    \end{cuenum}} \\ \hline
  \textbf{Postcondiciones} & \cuCelda{\begin{cureglas}
      \item \textbf{\mbox{POST-01}} Si el caso termina por el FN, el \textit{Usuario} tiene una \texttt{contrasena\_\allowbreak{}hash} y una \texttt{contrasena\_\allowbreak{}sal} nuevas, y la contraseña anterior ya no permite el acceso.
      \item \textbf{\mbox{POST-02}} Si el caso termina por el FN, la única \textit{Sesion} abierta de la cuenta es la que realizó el cambio, y ningún \textit{CodigoVerificacion} de la cuenta sigue pendiente.
      \item \textbf{\mbox{POST-03}} El \texttt{correo} y el \texttt{estado\_\allowbreak{}cuenta} no cambian en ningún flujo.
      \item \textbf{\mbox{POST-04}} Si el caso termina por el \mbox{FA-04} con bloqueo, ninguna \textit{Sesion} de la cuenta sigue abierta.
      \item \textbf{\mbox{POST-05}} En todos los casos que llegan al servidor, el intento quedó registrado en la \textit{BitacoraAcceso}.
    \end{cureglas}} \\ \hline
  \textbf{Reglas de negocio} & \cuCelda{\begin{cureglas}
      \item \textbf{\mbox{RN-01}} Cambiar la contraseña exige la contraseña actual, aunque haya una sesión abierta. La sesión prueba que alguien usa el dispositivo, no que sea el dueño de la cuenta.
      \item \textbf{\mbox{RN-02}} Las contraseñas no se almacenan, transmiten ni registran en claro.
      \item \textbf{\mbox{RN-03}} La contraseña nueva cumple la misma política que el alta (\mbox{RN-03} del \mbox{CU-02}).
      \item \textbf{\mbox{RN-04}} Cambiar la contraseña cierra todas las demás sesiones de la cuenta, pero conserva la que hizo el cambio, porque quien cambia la contraseña ya demostró conocer la anterior.
      \item \textbf{\mbox{RN-05}} Todo cambio de contraseña se avisa al correo de la cuenta. El aviso queda fuera de la transacción: que falle no deshace el cambio.
      \item \textbf{\mbox{RN-06}} La contraseña nueva debe ser distinta de la actual.
      \item \textbf{\mbox{RN-07}} Una contraseña actual incorrecta cuenta como intento fallido para el bloqueo escalonado, igual que en el inicio de sesión. De otro modo, este caso sería una forma de probar contraseñas sin límite desde una sesión robada.
      \item \textbf{\mbox{RN-08}} Los invitados no usan este caso: no tienen contraseña ni correo.
      \item \textbf{\mbox{RN-09}} El correo de la cuenta es el que se verificó en el registro (\mbox{CU-02}) y no puede cambiarse desde el juego (\mbox{D-28}). Es el canal que recibe los avisos de seguridad, como el del \mbox{RN-05}.
    \end{cureglas}} \\ \hline
  \textbf{Observación} & \cuCelda{\textit{Sobre por qué este caso solo cambia la contraseña.}\par
    La versión anterior de este caso cambiaba también el correo y activaba el segundo factor. Eran tres metas distintas con flujos distintos, y el nombre ``Cambiar contraseña o correo'' no decía cuál se perseguía. El cambio de correo se eliminó del sistema (\mbox{D-28}) y la activación del segundo factor pasó a su propio caso, el \mbox{CU-05}.} \\ \hline
  \textbf{Datos que toca} & \cuCelda{\begin{cudatos}
      \item \textbf{\textit{Sesion}:} lee por token; cierra las demás \textit{(pasos 6, 11)}
      \item \textbf{\textit{Usuario}:} lee \texttt{contrasena\_\allowbreak{}hash}, \texttt{contrasena\_\allowbreak{}sal}, \texttt{correo}, \texttt{idioma\_\allowbreak{}preferido} \textit{(pasos 7, 15)}
      \item \textbf{\textit{Usuario}:} escribe \texttt{contrasena\_\allowbreak{}hash}, \texttt{contrasena\_\allowbreak{}sal}, \texttt{intentos\_\allowbreak{}fallidos}, \texttt{fecha\_\allowbreak{}ultimo\_\allowbreak{}intento\_\allowbreak{}fallido}, \texttt{bloqueada\_\allowbreak{}hasta} \textit{(paso 10, \mbox{FA-04})}
      \item \textbf{\textit{CodigoVerificacion}:} invalida los pendientes \textit{(paso 12)}
      \item \textbf{\textit{BitacoraAcceso}:} inserta \textit{(paso 14, \mbox{FA-04})}
    \end{cudatos}} \\ \hline
\end{fichacu}
\clearpage
